% ===================================================================
%  TABELO N° 14 — La strado. — La komercisti.
%  Feuillets 91 a 97 du fac-simile = folios 89 a 95.
%  Correspondance : folio imprime = numero de feuillet - 2.
%
%  Les marques « %%K » ne sont pas des commentaires ordinaires : elles
%  portent la CLE D'ALIGNEMENT qui apparie, dans la page de lecture, ce
%  bloc-ci et son homologue francais. La cle se lit
%      t<tableau>-<numero d'alinea numerote>-<rang dans l'alinea>
%  et ne depend d'aucune pagination : les deux editions ne coupent pas
%  leurs pages aux memes endroits, mais elles numerotent les memes
%  alineas — c'est le seul point d'ancrage que les deux livrets
%  partagent, et il est de l'auteur, non du transcripteur.
% ===================================================================

% --- Feuillet 91 (folio 89 : ouverture de tableau) ------------------
\begin{VUpage}[91]{89}
%%K t14-apar-1 apar
\VUsaut{4.0mm}
\VUtitre{15.0pt}{60}{TABELO N\textsuperscript{o} 14}
\VUsaut{3.0mm}

%%K t14-tit-1 sub
\VUpk{11.4pt}{40}{La Strado. --- La Komercisti.}
\VUsaut{3.0mm}

%%K t14-01-1 p
\VUblancAlinea
1. --- Mea amiko Ioannes, qua habitas la\nl
ruro, venis pasar tri dii en mea familio. \VUgras{Til}\nl
lore il ne havabis l'okaziono vidar grand urbo;\nl
pro to ni uzas nia omna jorni promenante, e\nl
me explikas a lu to quon il ne konocas.

%%K t14-02-1 p
\VUblancAlinea
2. --- Unesme ni iris vizitar l'\VUgras{observatorio} \textsuperscript{(1)},\nl
qua tre interesis lu. Retrovenante tra la \VUgras{stra}\cc
\VUgras{do} \textsuperscript{(3)} ube esas la \VUgras{Turka balnerio} \textsuperscript{(2)}, ni renkon\cc
tris \VUgras{funero} \textsuperscript{(4)}. Avan la \VUgras{funeral konvoyo} mar\cc
chis la \VUgras{klerikaro,} preiranta la \VUgras{mortinto-ve}\cc
\VUgras{turo,} en qua esis pozita la \VUgras{sarko.} Mult amiki\nl
akompanis la \VUgras{trauranta} familio.

%%K t14-02-2 p
\VUblancAlinea
Pos la paso di la sequantaro, ni transiris la\nl
strado avan la granda \VUgras{birerio} \textsuperscript{(5)} en qua multa\nl
\VUgras{konsumanti} drinkis \VUgras{biro} sur la teraso, shir\cc
mate dal vitrizita \VUgras{marquizo} \textsuperscript{(6)}.

%%K t14-03-1 p
\VUblancAlinea
3. --- Ioannes esis astonegata da la \VUgras{skudo}\nl
pozita inter la magazino di la \VUgras{liquoristo} \textsuperscript{(7)} e\nl
ta di la \VUgras{muzik-vendisto} \textsuperscript{(10)}. Il questionis me\nl
pri olua signifiko; me respondis ke ol indikas\nl
la \VUgras{konsuleyo} \textsuperscript{(8)} Angla, e remarkigis da ilu la\nl
\VUgras{konsulo} \textsuperscript{(9)} qua esas sur la balkono.

%%K t14-tit-2 sub
\VUpk{10.2pt}{40}{Exploro al butiki.}
\VUsaut{3.0mm}

%%K t14-04-1 p
\VUblancAlinea
4. --- Kompreneble ni haltis avan l'estaleyo\nl
dil \VUgras{muzik-instrumenti,} \VUgras{harmoniumi,} \VUgras{orgeni} e\nl
\VUgras{piani}; ni regardis l'ilustrita kovrili di la par\cc
\parplein
\end{VUpage}

% --- Feuillet 92 (folio 90) -----------------------------------------
\begin{VUpage}[92]{90}
%%K t14-04-1 p suite
\VUcontinue
\VUgras{ticioni} e dil \VUgras{muzikaji} por piano, ed admiris la\nl
\VUgras{liro} ek orizita ligno, qua uzesas kom insigno.

%%K t14-05-1 p
\VUblancAlinea
5. --- La polva nubi facita dal \VUgras{balayisti} \textsuperscript{(11)} e\nl
la \VUgras{balayo-veturo} \textsuperscript{(12)} koaktis ni rapide pasar\nl
avan bela \VUgras{moblari} dil \VUgras{moblo-vendisto} \textsuperscript{(13)} ed\nl
avan l'estaleyo di \VUgras{furneli} e \VUgras{lampi} dil \VUgras{lad-(if-)}\nl
\VUgras{isto} \textsuperscript{(14)}.

%%K t14-06-1 p
\VUblancAlinea
6. --- Cetere, icaloke, ni koaktesis livar la\nl
trotuaro, nam ol esis inkombrita per paki,\nl
quin on prenis de \VUgras{transporto-veturo} \textsuperscript{(16)} por\nl
lokizar li en \VUgras{magazino} \textsuperscript{(15)} quan on jus loka\cc
cabis.

%%K t14-06-2 p
\VUblancAlinea
To impedis ni vidar la bela veturi dil \VUgras{vetu}\cc
\VUgras{rifisto} \textsuperscript{(17)}, ma ne impedis ni haltar por regar\cc
dar la \VUgras{pakigisto} \textsuperscript{(19)} facanta kesto por sua vi\cc
cino, la \VUgras{vendisto di cikli ed automobili} \textsuperscript{(18)}.\nl
Pose, pro ke ja esis kelke tarde, ni retrovenis\nl
a la hemo.

%%K t14-tit-3 sub
\VUpk{10.2pt}{40}{Promeno tra la urbo.}
\VUsaut{3.0mm}

%%K t14-07-1 p
\VUblancAlinea
7. --- Ye la morga dio, mea matro propozis\nl
a Ioannes ke on fotografez lu, pose el komisis\nl
me por akompanar lu che la \VUgras{fotografisto} \textsuperscript{(20)}.\nl
Pos la dejuno, ni iris aden la \VUgras{laboreyo} di l'ar\cc
tisto, che qua ni koaktesis sat longe vartar.\nl
Por okupar ni, ni regardis la \VUgras{portreti} \textsuperscript{(22)} pen\cc
danta an la muro.

%%K t14-07-2 p
\VUblancAlinea
Kande esis nia foyo, la laboro ne duris longe,\nl
e quik kande mea amiko finis pozar avan la\nl
\VUgras{fotografilo} \textsuperscript{(21)}, ni decensis.

%%K t14-08-1 p
\VUblancAlinea
8. --- Sur l'unesma etajo, ni vidis tra la\nl
pordo dil \VUgras{kuafisto} \textsuperscript{(23)}, olqua esas apertita, lua\nl
garsono qua tondis la hararo di kliento.
\end{VUpage}

% --- Feuillet 93 (folio 91) -----------------------------------------
\begin{VUpage}[93]{91}
%%K t14-08-2 p
\VUblancAlinea
Arivinte en la strado, ni pasis avan la \VUgras{tabak}\cc
\VUgras{vendeyo} \textsuperscript{(24)}. Ioannes esis tante amuzata da la\nl
reda \VUgras{lanterno} \textsuperscript{(25)} e la \VUgras{grosa pipo} qua esas uzata\nl
kom \VUgras{insigno} \textsuperscript{(26)}, ke il kolizionis kun du olda\nl
siori qui konversis \VUgras{sniflante tabako} \textsuperscript{(27)}.

%%K t14-tit-4 sub
\VUpk{10.2pt}{40}{La kukiferio.}
\VUsaut{3.0mm}

%%K t14-09-1 p
\VUblancAlinea
9. --- La omnalanda \VUgras{banko-bilieti} e \VUgras{monet}\cc
\VUgras{peci} expozita en la \VUgras{vetrino} dil \VUgras{kambiisto} \textsuperscript{(29)}\nl
kelkainstante atraktis nia atenco, ma pro ke\nl
esis la kloko repastetar, ni eniris che la \VUgras{ku}\cc
\VUgras{kifisto} \textsuperscript{(31)} sen haltar plu longe.

%%K t14-10-1 p
\VUblancAlinea
10. --- Vidante omna \VUgras{friandaji} quin kontenis\nl
la magazino, \VUgras{kuki, miel-kuki, bisquiti} ed om\cc
nasorta \VUgras{bonboni,} ni ne facile povis selektar.\nl
Fine Ioannes selektis la \VUgras{krem-kuki,} e me pre\cc
nis nur mikra \VUgras{ceriz-tarto,} nam la sukraji do\cc
\VUgras{lorigas mea denti,} e me tote ne deziras irar tro\nl
ofte vizitar la \VUgras{dentisto} \textsuperscript{(30)}.

%%K t14-tit-5 sub
\VUpk{10.2pt}{40}{Mashino-skribado.}
\VUsaut{3.0mm}

%%K t14-11-1 p
\VUblancAlinea
11. --- Por montrar a mea amiko \VUgras{skribo}\cc
\VUgras{mashino} pri qua il audabis parolar, ma quan\nl
il ne konocis, ni eniris la \VUgras{skribo-chambro} \textsuperscript{(32)}\nl
di mea \VUgras{kuzulo} \textsuperscript{(34)}. Ni trovis ilu diktanta sua\nl
letri a sua \VUgras{sekretario} \textsuperscript{(35)} qua esas \VUgras{stenogra}\cc
\VUgras{fisto}. Apene kande letro esis finita, \VUgras{mashin}\cc
\VUgras{skribisto} \textsuperscript{(33)} rikopiis olu per sua mashino. De\cc
zirante ne interruptar la okupesi di mea pa\cc
rento, ni restis che lu nur dum kelka instanti.

%%K t14-12-1 p
\VUblancAlinea
12. --- Sub la kontoro di mea kuzulo, habi\cc
tas granda \VUgras{horlojo-juvelisto} \textsuperscript{(37)} qua esas pluse
\parplein
\end{VUpage}

% --- Feuillet 94 (folio 92) -----------------------------------------
\begin{VUpage}[94]{92}
%%K t14-12-1 p suite
\VUcontinue
orlaboristo. Lua splendida magazino kontenas\nl
multa \VUgras{horloji, pendul-horloji, posh-horloji,}\nl
\VUgras{kateneti} e precoza \VUgras{juveli,} exemple \VUgras{perlo-kolia}\cc
\VUgras{ri,} ora \VUgras{braceleti, orel-ringi} e \VUgras{fingro-ringi} or\cc
nita per \VUgras{lapidi} e \VUgras{diamanti}. Un de la vetrini esas\nl
aparte rezervita por la \VUgras{oraji} e kontenas impor\cc
tanta kolektajo de arjenta \VUgras{manjilari}.

%%K t14-13-1 p
\VUblancAlinea
13. --- Turnante ye l'angulo di la strado, me\nl
remarkis ke on chanjabis la \VUgras{plako} \textsuperscript{(36)} pozita\nl
proxim la \VUgras{numero} di la domo. Me esis laute\nl
diconta mea remarko, kande la joyoza soni di\nl
\VUgras{la} militistal muziko esis audata. Ni kureskis\nl
renkontre a la regimento, sen reflektir ke \VUgras{lu}\nl
esis pasonta avan la magazini dil \VUgras{chapelisto} \textsuperscript{(39)}\nl
e dil \VUgras{gantisto} \textsuperscript{(40)}, e ke ni esabus tam bone\nl
lokizita por vidar olua defilo restante avan la\nl
vetrino dil \VUgras{optikisto} \textsuperscript{(38)}, tote provizita per
\parplein
\end{VUpage}

% --- Feuillet 95 (folio 93) -----------------------------------------
\begin{VUpage}[95]{93}
%%K t14-13-1 p suite
\VUcontinue
\VUgras{naz-binokli, orel-binokli} e \VUgras{lorneti}.

%%K t14-tit-6 sub
\VUpk{10.2pt}{40}{Marcho-defilo.}
\VUsaut{3.0mm}

%%K t14-14-1 p
\VUblancAlinea
14. --- Avan la \VUgras{regimento} \textsuperscript{(41)} marchis la \VUgras{tam}\cc
\VUgras{burestro} \textsuperscript{(42)} sequata da la \VUgras{tamburisti} \textsuperscript{(43)}, la\nl
\VUgras{klarionisti} \textsuperscript{(44)} e la tota \VUgras{muzikistaro}. La gene\cc
\VUgras{ralo} \textsuperscript{(45)} qua esis sur la placo kun sua adjutanto,\nl
esis haltinta proxim la \VUgras{aquo-spricilo} \textsuperscript{(49)} di la\nl
\VUgras{fonteno} \textsuperscript{(48)}, por asistar la \VUgras{defilo}.

%%K t14-14-2 p
\VUblancAlinea
\VUgras{La stab-oficiri} \textsuperscript{(46)}, la \VUgras{kolonelo} e la \VUgras{lietnanto}\cc
\VUgras{kolonelo} salutis lu per espado. \VUgras{La mayori} esis\nl
avan sua \VUgras{batalioni} \textsuperscript{(47)} e singla \VUgras{kapitano} ko\cc
mandis sua \VUgras{kompanio,} dum ke la \VUgras{lietnanti,}\nl
\VUgras{sublietnanti} e \VUgras{suboficiri} okupis sua kustumata\nl
plaso apud sua viri.

%%K t14-15-1 p
\VUblancAlinea
15. --- Multa pasanti grupeskabis sur \VUgras{la}\nl
\VUgras{voyo-kruco} \textsuperscript{(50)}; la komercisti, la \VUgras{parapluvis}\cc
\VUgras{to} \textsuperscript{(51)}, la \VUgras{tintisto} \textsuperscript{(53)}, l'\VUgras{armifisto} \textsuperscript{(54)} ekiris por\nl
vidar la \VUgras{soldati}. Omna vehili, \VUgras{fiakri} \textsuperscript{(52)}, \VUgras{mas}\cc
\VUgras{tro-veturi} \textsuperscript{(57)} ed anke l'\VUgras{aroz-fudro} \textsuperscript{(56)}, garis\nl
su alonge la trotuaro.

%%K t14-tit-7 sub
\VUpk{10.2pt}{40}{Ceni an la strado.}
\VUsaut{3.0mm}

%%K t14-15-2 p
\VUblancAlinea
\VUgras{Komerco-voyajisto} \textsuperscript{(55)} portanta per manuo\nl
sua \VUgras{specimenal buxi,} rapide transiris la strado,\nl
e la personi sidanta sur l'\VUgras{emperialo} \textsuperscript{(59)} dil \VUgras{om}\cc
\VUgras{nibuso} \textsuperscript{(58)} riturnis su por juar la vidindajo.\nl
Kompatinda \VUgras{kantero vagera} \textsuperscript{(60)} kun sua \VUgras{turn}\cc
\VUgras{orgeno} \textsuperscript{(61)} mustis interruptar sua \VUgras{kansono,}\nl
nam la bruisado dil tamburi kovris lua voco.

%%K t14-16-1 p
\VUblancAlinea
16. --- Kande la regimento arivis avan la\nl
\VUgras{stof-komerceyo} \textsuperscript{(64)} la \VUgras{sutistini} \textsuperscript{(63)} e la \VUgras{talio}\cc
\VUgras{ri} \textsuperscript{(62)} livis sua \VUgras{suto-mashini} e sua laboro por\nl
regardar tra la fenestro. \VUgras{La komercisto} \textsuperscript{(65)}, qua\nl
jus pozabis nova \VUgras{kostumi} \textsuperscript{(66)} an sua \VUgras{estaleyo,}\nl
devis pozigar interne la \VUgras{drapaji} \textsuperscript{(67)} e la \VUgras{telaji}\nl
instalita sur la trotuaro, proxim la \VUgras{kloak-aper}\cc
\VUgras{turo} \textsuperscript{(68)}, time ke la turbo renversos li, mem\nl
olda \VUgras{kolportero} \textsuperscript{(69)}, de qua pordistino mar\cc
chandis kalzi, esis koaktata enirar koridoro\nl
por finar sua vendo.

%%K t14-16-2 p
\VUblancAlinea
Ni akompanis la regimento til la kazerno, \VUgras{e,}\nl
tre kontenta pri nia jornedo, ni retrovenis ye\nl
la horo dil dineo.

%%K t14-tit-8 sub
\VUpk{10.2pt}{40}{Diversa komerci e profesioni.}
\VUsaut{3.0mm}

%%K t14-17-1 p
\VUblancAlinea
17. --- En la sequanta dio, mea patrulo pro\cc
pozis a ni vizitar l'imprimerio dil iba jurnalo.\nl
Ni entuziasme aceptis.
\end{VUpage}

% --- Feuillet 96 (folio 94) -----------------------------------------
\begin{VUpage}[96]{94}
%%K t14-17-2 p
\VUblancAlinea
L'\VUgras{imprimisto} esas anke \VUgras{libristo} \textsuperscript{(70)} \VUgras{(=libro}\cc
\VUgras{vendisto),} \VUgras{editisto,} \VUgras{papervendisto} e \VUgras{bindisto}.\nl
Il instalis sur l'unesma etajo la \VUgras{redakteyo} \textsuperscript{(71)}\nl
dil jurnalo, ed anke la \VUgras{grabeyo,} en qua laboras\nl
la \VUgras{litografisti} \textsuperscript{(72)} e la \VUgras{kupro-grabisti} \textsuperscript{(73)}, fine\nl
la \VUgras{komposto-laboreyo,} en qua esas la \VUgras{imprim}\cc
\VUgras{laboristi} \textsuperscript{(74)}. L'\VUgras{imprimeyo} \textsuperscript{(75)} propre dicata, la\nl
\VUgras{imprim-presili} e la diversa mashini esas sur\nl
la ter-etajo.

%%K t14-tit-9 sub
\VUpk{10.2pt}{40}{Ioannes facas tre multa questioni.}
\VUsaut{3.0mm}

%%K t14-18-1 p
\VUblancAlinea
18. --- Ekirante l'imprimerio, Ioannes tre\nl
astonita haltis avan mikra vitrizita buxo po\cc
zita sur kolono, kontre la \VUgras{mercereyo} \textsuperscript{(77)} e ques\cc
tionis por quo ol uzesas. Mea patrulo explikis\nl
ad ilu, ke ol esas \VUgras{incendio-avertilo} \textsuperscript{(76)}. Cetere\nl
omno, en la urbo, esis por mea amiko, asto\cc
nivo, de la simpla \VUgras{ston-fonteno} \textsuperscript{(78)} e la lejera\nl
\VUgras{porto-triciklo} \textsuperscript{(80)} duktata da \VUgras{grumo} \textsuperscript{(79)} kun li\cc
vreo, til la \VUgras{posto-veturo} \textsuperscript{(81)}.

%%K t14-19-1 p
\VUblancAlinea
19. --- Dum ke mea patrulo kelkinstante livis\nl
ni por parolar kun l'\VUgras{imposto-kolektisto} \textsuperscript{(83)},\nl
qua haltabis por konversar kun \VUgras{advokato} \textsuperscript{(82)} e\nl
\VUgras{notario} \textsuperscript{(84)}, ni amuzis ni sequante per la okuli\nl
la cirkulado sur la strado. La maxim diversa\nl
personi pasis avan ni : \VUgras{kukista garsono} \textsuperscript{(85)},\nl
portanta en korbo \VUgras{pasteto} splendida, venis dop\nl
\VUgras{vesto-vendisto} \textsuperscript{(86)}; \VUgras{lantern-acendisto} \textsuperscript{(87)} esis\nl
konversanta kun \VUgras{gasisto} \textsuperscript{(88)}, \VUgras{regulierino} \textsuperscript{(89)} te\cc
nanta manue orfanino esis retroiranta aden\nl
sua kuvento.
\end{VUpage}

% --- Feuillet 97 (folio 95) -----------------------------------------
\begin{VUpage}[97]{95}
%%K t14-tit-10 sub
\VUpk{10.2pt}{40}{Dum nia retroveno adheme.}
\VUsaut{3.0mm}

%%K t14-20-1 p
\VUblancAlinea
20. --- Ye l'instanto kande mea patrulo ra\cc
juntis ni, Ioannes ridegis, vidante la sordida\nl
vizajo e la stranja vestizuracho di du \VUgras{kameno}\cc
\VUgras{skrapisti} \textsuperscript{(91)} tote fuliginoza.

%%K t14-21-1 p
\VUblancAlinea
21. --- Me volis komprar de la \VUgras{buketistino} \textsuperscript{(92)}\nl
planto por mea matro, ma mea patrulo remar\cc
kigis da me, ke ol esus tre pezoza por kunpor\cc
tar e ke esus preferinda aquirar \VUgras{oranji}. Ni\nl
proximeskis a la \VUgras{vendistino} \textsuperscript{(94)} e kande la\nl
\VUgras{edukistino} \textsuperscript{(93)}, qua esis selektanta por sua edu\cc
katulo, finis sua kompro, ni igis servar ni.

%%K t14-22-1 p
\VUblancAlinea
22. --- Retrovenante a la hemo, ni renkontris\nl
plura konocati : olda amikino di mea familio,\nl
elqua apogis la brakio sur sua \VUgras{kompana siori}\cc
\VUgras{no} \textsuperscript{(95)}, l'\VUgras{injenioro} \textsuperscript{(96)} pri la ponti e voyi, ed\nl
un de mea kuzini kun lua \VUgras{infantala gardis}\cc
\VUgras{tino} \textsuperscript{(98)}.

%%K t14-22-2 p
\VUblancAlinea
Pose ni renkontris la \VUgras{chapelistino} \textsuperscript{(97)} di mea\nl
matro e du \VUgras{monaki} \textsuperscript{(99)} stranjera a la urbo, qui\nl
adiris ni por questionar mea patrulo pri la\nl
voyo a la staciono.

\VUsaut{6.0mm}
\VUfilet{22mm}
\end{VUpage}
